\textbf{Phần 1. Trắc nghiệm lựa chọn}
\setcounterpageref{ex}{0}

\begin{ex} %[2D10V48]
Số gồm $10$ chục que tính gộp lại được gọi là số nào?
\choice
{Mười chục}
{Một nghìn}
{\True Một trăm}
{Mười đơn vị}
\loigiai{
$10$ chục $= 100$ (một trăm).
\textbf{Đáp án của câu hỏi là C}
}
\end{ex}

\begin{ex} %[2D10V48]
Số gồm $10$ trăm đơn vị gộp lại được gọi là số nào?
\choice
{\True Một nghìn}
{Một trăm}
{Mười chục}
{Một vạn}
\loigiai{
$10$ trăm $= 1 000$ (một nghìn).
\textbf{Đáp án của câu hỏi là A}
}
\end{ex}

\begin{ex} %[2D10V48]
Một người xếp các túi bi, mỗi túi có $10$ viên bi. Có $10$ túi như thế đóng trong một hộp. Hỏi một hộp có tất cả bao nhiêu viên bi?
\choice
{$20$ viên}
{$50$ viên}
{$1 000$ viên}
{\True $100$ viên}
\loigiai{
$10$ túi, mỗi túi 10 viên là $10$ chục viên $= 100$ viên bi.
\textbf{Đáp án của câu hỏi là D}
}
\end{ex}

\begin{ex} %[2D10V48]
Một kho hàng có $10$ thùng sách, mỗi thùng chứa đúng $100$ cuốn sách. Hỏi trong kho có tất cả bao nhiêu cuốn sách?
\choice
{$100$ cuốn}
{\True $1 000$ cuốn}
{$500$ cuốn}
{$200$ cuốn}
\loigiai{
$10$ thùng, mỗi thùng 100 cuốn là $10$ trăm cuốn $= 1 000$ cuốn sách.
\textbf{Đáp án của câu hỏi là B}
}
\end{ex}

\begin{ex} %[2D10V48]
Hiệu giữa số tự nhiên bé nhất có bốn chữ số và số tự nhiên lớn nhất có ba chữ số là
\choice
{$10$}
{$100$}
{\True $1$}
{$0$}
\loigiai{
Số bé nhất có bốn chữ số là $1 000$, số lớn nhất có ba chữ số là $999$. Ta có: $1 000 - 999 = 1$.
\textbf{Đáp án của câu hỏi là C}
}
\end{ex}

\begin{ex} %[2D10V48]
Trong số $1 000$, chữ số $1$ nằm ở hàng nào?
\choice
{\True Hàng nghìn}
{Hàng đơn vị}
{Hàng chục}
{Hàng trăm}
\loigiai{
Số $1 000$ gồm $1$ nghìn, $0$ trăm, $0$ chục, $0$ đơn vị.
\textbf{Đáp án của câu hỏi là A}
}
\end{ex}

\begin{ex} %[2D10V48]
Có bao nhiêu chục que tính để gộp lại thành đúng $4$ trăm que tính?
\choice
{$4$ chục}
{$14$ chục}
{$400$ chục}
{\True $40$ chục}
\loigiai{
$4$ trăm $= 40$ chục.
\textbf{Đáp án của câu hỏi là D}
}
\end{ex}

\begin{ex} %[2D10V48]
Bác Nam đóng gói cam vào $7$ thùng to, mỗi thùng $100$ quả và $3$ túi nhỏ, mỗi túi $10$ quả. Bác Nam có tất cả bao nhiêu quả cam?
\choice
{$703$ quả}
{\True $730$ quả}
{$73$ quả}
{$370$ quả}
\loigiai{
$7$ trăm và $3$ chục quả cam là $730$ quả cam.
\textbf{Đáp án của câu hỏi là B}
}
\end{ex}

\begin{ex} %[2D10V48]
Một xưởng sản xuất đóng được $1 000$ chiếc quạt. Người ta chở đi phân phối bằng các xe tải nhỏ, mỗi xe chở được $100$ chiếc. Cần bao nhiêu chuyến xe để chở hết số quạt đó?
\choice
{$5$ chuyến}
{$100$ chuyến}
{\True $10$ chuyến}
{$8$ chuyến}
\loigiai{
$1 000$ gồm $10$ trăm, do đó cần $10$ chuyến xe.
\textbf{Đáp án của câu hỏi là C}
}
\end{ex}

\begin{ex} %[2D10V48]
Số gồm $9$ trăm và $10$ chục là số nào?
\choice
{\True $1 000$}
{$910$}
{$901$}
{$900$}
\loigiai{
$10$ chục $= 1$ trăm. $9$ trăm thêm $1$ trăm là $10$ trăm $= 1 000$.
\textbf{Đáp án của câu hỏi là A}
}
\end{ex}

\begin{ex} %[2D10V48]
Một cửa hàng bán hoa kết các bó hoa hồng, mỗi bó có $10$ bông. Để có đủ $300$ bông hoa, cửa hàng cần chuẩn bị bao nhiêu bó hoa như vậy?
\choice
{$3$ bó}
{$300$ bó}
{$33$ bó}
{\True $30$ bó}
\loigiai{
$300$ bông gồm $30$ chục bông, vậy cần $30$ bó hoa.
\textbf{Đáp án của câu hỏi là D}
}
\end{ex}

\begin{ex} %[2D10V48]
Số $1 000$ có giá trị gấp bao nhiêu lần số $100$?
\choice
{$5$ lần}
{\True $10$ lần}
{$100$ lần}
{$2$ lần}
\loigiai{
$1 000 = 100 \times 10$, nên số $1 000$ gấp $10$ lần số $100$.
\textbf{Đáp án của câu hỏi là B}
}
\end{ex}

\textbf{Phần 2. Trắc nghiệm đúng sai}
\setcounterpageref{ex}{0}

\begin{ex} %[2D10V48]
Xét cấu tạo và mối quan hệ giữa các hàng đơn vị, chục, trăm, nghìn:
\choiceTF
{\True Mười chục đơn vị gộp lại bằng một trăm.}
{Mười trăm đơn vị gộp lại bằng mười nghìn.}
{\True Số liền trước của số 1 000 là số 999.}
{\True Một nghìn gồm một trăm chục.}
\loigiai{
\begin{itemize}
    \item Ý a) Đúng: 10 chục = 100 (một trăm).
    \item Ý b) Sai: 10 trăm = 1 nghìn, không phải 10 nghìn.
    \item Ý c) Đúng: 999 + 1 = 1 000.
    \item Ý d) Đúng: 100 chục = 100 x 10 = 1 000.
\end{itemize}
\begin{itemize}
    \item a) \textbf{Đúng}.
    \item b) \textbf{Sai}.
    \item c) \textbf{Đúng}.
    \item d) \textbf{Đúng}.
\end{itemize}
}
\end{ex}

\begin{ex} %[2D10V48]
Một xưởng bánh có 8 thùng bánh to (mỗi thùng 100 cái) và 20 hộp bánh nhỏ (mỗi hộp 10 cái):
\choiceTF
{\True Số bánh trong 8 thùng to là 800 cái.}
{\True Số bánh trong 20 hộp nhỏ là 200 cái.}
{\True Tổng số bánh xưởng có là 1 000 cái.}
{Xưởng cần thêm 100 cái bánh nữa để đủ 1 000 cái.}
\loigiai{
\begin{itemize}
    \item Ý a) Đúng: 8 x 100 = 800 cái.
    \item Ý b) Đúng: 20 x 10 = 200 cái.
    \item Ý c) Đúng: 800 + 200 = 1 000 cái bánh.
    \item Ý d) Sai: Xưởng đã có đủ 1 000 cái rồi.
\end{itemize}
\begin{itemize}
    \item a) \textbf{Đúng}.
    \item b) \textbf{Đúng}.
    \item c) \textbf{Đúng}.
    \item d) \textbf{Sai}.
\end{itemize}
}
\end{ex}

\textbf{Phần 3. Trả lời ngắn}
\setcounterpageref{ex}{0}

\begin{ex} %[2D10V48]
Tìm số liền sau của số lớn nhất có ba chữ số.
\shortans{1 000}
\loigiai{
\textbf{Đáp số:} $1 000$\\
Số lớn nhất có 3 chữ số là 999. Số liền sau là 1 000.
}
\end{ex}

\begin{ex} %[2D10V48]
Một cửa hàng có 6 hộp bút lớn (mỗi hộp 100 chiếc) và 4 hộp bút nhỏ (mỗi hộp 10 chiếc). Cửa hàng có tất cả bao nhiêu chiếc bút?
\shortans{640}
\loigiai{
\textbf{Đáp số:} $640$\\
6 hộp lớn có 600 chiếc, 4 hộp nhỏ có 40 chiếc. Tổng là $600 + 40 = 640$ chiếc.
}
\end{ex}

\begin{ex} %[2D10V48]
Cần thêm bao nhiêu chục que tính vào 8 trăm que tính để được đúng 1 000 que tính?
\shortans{20}
\loigiai{
\textbf{Đáp số:} $20$\\
1 000 que tính = 10 trăm. Cần thêm $10 - 8 = 2$ trăm. 2 trăm = 20 chục.
}
\end{ex}

\begin{ex} %[2D10V48]
Có bao nhiêu chữ số 0 trong cách viết số một nghìn?
\shortans{3}
\loigiai{
\textbf{Đáp số:} $3$\\
Số 1 000 có 3 chữ số 0.
}
\end{ex}

\begin{ex} %[2D10V48]
Một sọt cam có 100 quả. Hỏi 7 sọt cam như thế có bao nhiêu quả cam?
\shortans{700}
\loigiai{
\textbf{Đáp số:} $700$\\
7 x 100 = 700 quả cam.
}
\end{ex}

\begin{ex} %[2D10V48]
Tìm số tự nhiên tròn chục sao cho khi lấy số đó nhân với 10 thì được 500.
\shortans{50}
\loigiai{
\textbf{Đáp số:} $50$\\
Số đó là: $500 : 10 = 50$.
}
\end{ex}
